\section{讨论、局限与结论}
\label{sec:discussion}
\subsection{有效性边界}
本实验是单变量、固定长度、确定性生成的评估场景。异常标签由生成协议规定，不来自独立专家标注。B 和 C 使用真实基线，因此它们在本样例中的表现不能直接外推到漂移、缺失、多变量或未知周期环境。更不能把这些合成指标当作设备故障识别的真实准确率。

表格中的六位数据和四位指标是显示精度，不代表测量精度；本样例没有真实测量误差模型。虽然附录包含 120 条记录，但它们不是 120 次独立随机试验，不能据此把标准误简单缩小为原来的 $1/\sqrt{120}$。

\subsection{文档系统的验收价值}
从软件角度看，成功生成 PDF 只是基础要求。还应检查：章节和目录编号是否一致，正文引用是否显示为问号，参考文献是否只有已引用的三条，PNG 是否清晰，跨页表格是否重复表头，以及修改数值后能否在下一次预览中看到变化。

故意删除图片、移除一个文献键或写错公式，可以作为错误定位测试，但这些破坏性场景应在副本中进行。默认交付文件保持可成功编译，不混入未闭合环境或不存在的资源。同步定位依赖编译器产生的 SyncTeX 文件，不能仅凭 PDF 页数推断其功能正常。

\subsection{结论}
本文给出了一套完整的中文多文件论文样例，将确定性数据生成、数学定义、三线表、长表格、图片和三条正文引用组合为可复现编译输入。样例的核心价值在于可核对、可重新生成且明确区分合成测试与真实研究。未来可以在独立副本中扩展长文档、大图片和复杂文献样式压力测试，但不应把额外宏包或外部工具悄悄加入当前离线依赖链。

\section*{数据声明与致谢}
全部数值、图像及内部测试文献由本样例生成或编写，不包含真实个人数据。作者署名为测试工作组，不对应真实学术机构。感谢本地编辑器提供编译验证环境；该表述不意味着任何第三方对实验结论背书。
